\subsubsection{数据链、算法原理与完整步骤}

数据工作从数据字典和泄漏边界开始。缺失值按形成机制选择均值/中位数、插值、
KNN 填补或删除；异常值用 $3\sigma$、箱线图和 Isolation Forest 交叉检查，
不能把业务极值自动删除；重复记录先定义主键和时间容差。尺度处理比较
Z-score 标准化与 Min--Max 归一化，类别变量按名义/有序属性选择 one-hot 或
label 编码，连续变量只有在解释或规则需要时才按等宽/等频分箱。样本筛选使用
Monte Carlo 随机抽样时，须固定种子，记录抽样分布并核对覆盖率；它与组员1
的不确定性仿真共享随机数配置，但不把“抽取样本”与“传播不确定性”混为同一任务。

探索阶段采用 Pearson、Spearman、Kendall 相关，分别对应线性、单调秩和序关系；
聚类比较 K-Means、DBSCAN、层次聚类和谱聚类。空间或时间缺口比较 Lagrange、
三次样条和 Kriging 插值。PCA 用于线性方差压缩，t-SNE 主要用于局部结构可视化，
UMAP 兼顾局部与部分全局结构；后两者的二维图不直接充当预测证据。

评价决策链覆盖 AHP、模糊综合评价、模糊 AHP、灰色关联、PCA、因子分析、熵权、
CRITIC、变异系数、TOPSIS、VIKOR、灰色评价和物元可拓，并训练熵权--AHP、
CRITIC--TOPSIS 等组合。主观权重必须给判断矩阵和一致性检验，客观权重必须解释
数据离散或冲突如何影响权重；TOPSIS 的距离、VIKOR 的折中和灰色关联的序列相似
并非同一含义，不能只因都输出排序而互换。

回归覆盖多元线性、Ridge、Lasso、SVR、随机森林、XGBoost 和 CatBoost；分类覆盖
Logistic、朴素 Bayes、SVM、决策树、随机森林、AdaBoost 和 CNN。时间序列覆盖
ARIMA、SARIMA、指数平滑、VAR、GM(1,1)、灰色--Markov、Prophet、XGBoost、
LightGBM、BP、RNN、LSTM、TCN 和 GRU。线性/广义线性模型保留可解释基线，树模型
处理非线性与交互，深度网络仅在样本规模、序列长度和验证方案足以支撑时启用。
思维导图中的“Logistic 预测模型”归入二分类模型，不当作连续响应回归；长期预测是
外推任务，必须报告预测区间、结构漂移风险和超出训练区间的限制。

完整步骤为：
\begin{enumerate}
  \item 建立字段、单位、主键、时间戳、缺失和异常审计表；
  \item 在任何填补、标准化、降维或特征选择前划分训练/验证边界；
  \item 完成探索图表与业务一致性核验，建立简单均值/线性/持续性基线；
  \item 按任务选择评价、回归、分类、聚类或时序模型，并列出近邻方法边界；
  \item 将预处理与模型封装在同一 pipeline，采用 K 折、留出或滚动窗口验证；
  \item 回归报告 MAE、MSE、RMSE、MAPE，分类报告混淆矩阵及类别不平衡指标；
  \item 做残差、校准、特征稳定性、阈值和数据扰动分析，再向决策模块输出。
\end{enumerate}

